\clearpage\thispagestyle{empty}
\pdfbookmark[0]{Quellen zur Kometenanzeige von 1814}{anzeige1814-sources}
\begin{center}\textbf{Quellen zur Kometenanzeige von 1814}\end{center}
\noindent Diese KI-Übertragung des deutschen Originals in heutiges Deutsch folgt dem Text in Gauß’ \textit{Werken}, Band~VI (Göttingen, 1874), S.~65--67. Die Anzeige erschien in den \textit{Göttingischen gelehrten Anzeigen}, 3.~Stück, 3.~Januar 1814, S.~25--28. \href{https://www.digitale-sammlungen.de/de/view/bsb10538685?page=31}{Erstdruck in der Bayerischen Staatsbibliothek}. \href{https://gdz.sub.uni-goettingen.de/id/PPN236007467}{Historischer deutscher Text in den Werken}.
\par\smallskip\noindent In den \textit{Werken} reichen die Beobachtungen an der Sternwarte vom 9. bis zum 25.~April. Im Erstdruck reichen sie vom 7. bis zum 25.~April. In beiden Fassungen verwendet das abschließende Beispiel Beobachtungen vom 7., 14. und 21.~April.
\par\smallskip\noindent Übertragung und Satz wurden mit KI erstellt.
\clearpage
